\chapter*{Editor's note: status at assembly}
\addcontentsline{toc}{chapter}{Editor's note: status at assembly}
The five chapters were written independently, each from the same fact sheet (\texttt{FACTS.md}), and combined
afterwards. This note records what changed while they were being written.

\paragraph{D4: open versus closed disks \tierA.} The theorem of Section~\ref{ch2:sec:d4} (open-disk density and lower
energy coincide with the closed-disk ones) is kernel-checked in \texttt{formal\_cert/Triangular/OpenDisk.lean} for
density one and for the lower energy. The density statement at other densities $\rho\ne1$ is \tierC{}: the same
argument applies, but it is not formalised. The Lean theorems are:
\begin{itemize}
\item \lean{energy_eq_openEnergy}, \lean{densityOne_iff_open} and \lean{universal_energy_minimum_open}, which is
upstream's theorem verbatim with open disks;
\item \lean{A_densityOneOpen}, the non-vacuity check.
\end{itemize}
It was preregistered before the proof was written. An independent verification run accepted it with corrections. That
run found that the two identities need no local-finiteness hypothesis at all, and kernel-checked this in
\texttt{OpenDiskAny.lean}. It also added a control showing that the identity is not definitional: for a single point on
the unit circle, the open and closed counts differ at $R=1$. Every footprint is within \{propext, Classical.choice,
Quot.sound\}.

\emph{Scope:} this settles the open/closed question for upstream's observables only (centred disks, $N/(\pi R^2)$
density, $\liminf$ energy in $[0,\infty]$). It says nothing about other window shapes, $\limsup$ energies or physical
boundary conditions.

\paragraph{Jellium (\tierX).} Status: \JelliumMainStatus.

\paragraph{Erratum carried from the published paper (found while writing Chapter~3).} The numerical study uses a
shifted-force truncation of the Yukawa potential at $r_c=8$. The quantity reported in the paper as the truncation size
is $V(r_c)$, which is \emph{not} the change in the energy per particle. Summing the unshifted potential out to $r=60$
changes the absolute energies by about $9.5\times10^{-6}$ at $\kappa=2$ and $1.8\times10^{-2}$ at $\kappa=1$; at
$\kappa=1$ that exceeds $e_{\rm sq}-e_{\rm tri}$. The square--triangular gap itself is almost unchanged, and every
comparison in the study uses the same truncated potential, so the P1 observation stands for that potential. But the
truncated potential is not completely monotone, so P1 concerns an approximation of the Yukawa interaction, not the
interaction the theorem covers. The next version of the paper will state this.

\paragraph{Scripts.} The Python scripts quoted in Chapters~3--5 are in \texttt{code/}. Chapter~1 prints its short
checks inline. The scratch checks quoted in Chapter~3's footnotes are not shipped.
